%----------------------------------------
% IT IS RECOMMENDED TO USE AUTOLATEX FOR
% COMPILING THIS DOCUMENT.
% http://www.arakhne.org/autolatex
%----------------------------------------

\documentclass[english,nodocumentinfo,nopubpage]{utbmciadreport}

% The TeX code is entering with UTF8
% character encoding (Linux and MacOS standards)
\usepackage[utf8]{inputenc}
\usepackage{fancyhdr}
\usepackage{../common/sarl-colorized-listing}


\graphicspath{{./}}

\declaredocument{Introduction to SARL}{Tutorial Session \#3}{UTBM-INFO-AI51-TU3}

\addauthorvalidator*[St\'ephane Galland]{St{\'e}phane}{Galland}{Teacher}

\updateversion{1.0}{\makedate{20}{03}{2023}}{First release on Github}{\upmpublic}

\begin{document}

\section{Goal of this Tutorial Session}

The goal of this tutorial session is to introduce the basics concepts of agent-oriented programming with SARL.
The general goal of the tutorial is to remind several key points related to the coding with SARL.

\section{Work to be Done during the Tutorial Session}

The following sections describe the work to be done during this tutorial session. It is recommended to work 20 minutes max on each exercice.

\subsection{Exercice 1 - Events and Basic Lifecycle}

The goal of this exercice is to remind the base elements related to the events and agent lifecycle in SARL.

\begin{enumerate}[a)]
\item Define an event with the type name \code{Presentation}\ that represents a message of presentation from one agent to another agent.
\item Define an agent with the type name \code{Agent1}\ that has the following behavior:
	\begin{itemize}
	\item When it is starting, it is displaying a log message that is indicating it is alive.
	\item When the agent is receiving an event occurrence of type \code{Presentation}, it is increasing an internal counter that is counting the number of presentations received by the agent.
	\item When the agent is destroying itself, it has to log the total number of presentations that were received during its live.
	\end{itemize}
\end{enumerate}


\subsection{Exercice 2 - Event Communication}

The goal of this exercice is to remind the base of the agent communication in SARL.

\begin{enumerate}[a)]
\item Define an agent with the type name \code{Agent2}\ that has the following behavior:
	\begin{itemize}
	\item When it is starting, it is emitting an event of type \code{Presentation}\ to notify all the other agents in the system that this new agent is alive.
	\end{itemize}
\item Create an agent with the type name \code{Agent3} that is spawning an agent of type \code{Agent1}\ and \code{Agent2}.
\item Kill all the agents for stopping the application.
\end{enumerate}

\subsection{Exercice 3 - AgentSpawned}

The goal of this exercice is to introduce the predefined events in SARL, that are dedicated to the agent lifecycle.

\begin{enumerate}[a)]
\item Define an agent with the type name \code{Agent4}\ that has the following behavior:
	\begin{itemize}
	\item After starting, wait 1 minute and kill itself.
	\end{itemize}
\item Define an agent with the type name \code{Agent5}\ that has the following behavior:
	\begin{itemize}
	\item Wait for the spawning of an agent and display the ID of the spawned agent.
	\end{itemize}
\item Update the code of the agent with the type name \code{Agent5}\ in order to:
	\begin{itemize}
	\item Spawn 5 agents of type \code{Agent4}.
	\item Wait for the spawning of the 5 agents.
	\item Kill all the agents \Emph{after} all the agents were spawned.
	\end{itemize}
\end{enumerate}

\subsection{Exercice 4 - Capacity and Skill}

The goal of this exercice is to have multiple agents that are implementing in different ways a know-how. This know-how is the capacity to save data into a text-based file format.

\begin{enumerate}[a)]
\item Define a capacity with the type name \code{DataSave}\ that provides the following function for saving data: \\
	\code{def saveData(data : List<?>) : void}
\item Define an implementation (named \code{JsonDataSave}) of the capacity \code{DataSave}\ that is saving the data with the JSON format into the file ``data.json''.
\item Define an implementation (named \code{XmlDataSave}) of the capacity \code{DataSave}\ that is saving the data with the XML format into the file ``data.xml''.
\item Create an agent with the type name \code{Agent6}\ that has the following behavior:
	\begin{itemize}
	\item Every 1 minute, compute randomly an list of integer values, and save it using the JSON format.
	\end{itemize}
\item Create an agent with the type name \code{Agent7}\ that has the following behavior:
	\begin{itemize}
	\item Every 1 minute, compute randomly an list of integer values, and save it using the XML format.
	\end{itemize}
\item Create an agent with the type name \code{Agent8}\ that has the following behavior:
	\begin{itemize}
	\item Spawn an agent of type \code{Agent5}.
	\item Spawn an agent of type \code{Agent7}.
	\end{itemize}
\end{enumerate}

\subsection{Exercice 5 - Dynamic Skill}

The goal of this exercice is to have an agent that is implementing in different ways a know-how. This know-how is the capacity to save data into a text-based file format.

\begin{enumerate}[a)]
\item Create an agent with the type name \code{Agent9}\ that has the following behavior:
	\begin{itemize}
	\item Every 1 minute, compute randomly an list of integer values, select the skill randomly between the JSON and XML savings, and save the random data using the selected skill.
	\end{itemize}
\item Create an agent with the type name \code{Agent10}\ that has the following behavior:
	\begin{itemize}
	\item Spawn an agent of type \code{Agent9}.
	\end{itemize}
\end{enumerate}

\subsection{Exercice 6 - Parallel Execution}

The goal of this exercice is to introduce the parallel execution of the agents.

\begin{enumerate}[a)]
\item Create an event with the type name \code{MyEvent}.
\item Create an agent with the type name \code{Agent11}\ that has the following behavior:
	\begin{itemize}
	\item One event handler that is executed when receiving an event occurrence of type \code{MyEvent}, and displaying in the logs the message ``First handler''.
	\item A second event handler that is executed when receiving an event occurrence of type \code{MyEvent}, and displaying in the logs the message ``Second handler''.
	\end{itemize}
\item Create an agent with the type name \code{Agent12}\ that has the following behavior:
	\begin{itemize}
	\item Spawn an agent of type \code{Agent11}.
	\item Every 5 seconds, emit an event of type \code{MyEvent}. The agent of type \code{Agent12}\ should display the two messages.
	\end{itemize}
\end{enumerate}


\subsection{Exercice 7 - Event handler guard}

The goal of this exercice is to introduce the parallel execution of the agents.

\begin{enumerate}[a)]
\item Create an agent with the type name \code{Agent13}\ that has the following behavior:
	\begin{itemize}
	\item One event handler that is executed when receiving an event occurrence of type \code{MyEvent}, and displaying in the logs the message ``First handler''.
	\item A second event handler that is executed when receiving an event occurrence of type \code{MyEvent}, and displaying in the logs the message ``Second handler''.
	\item When an event of type \code{MyEvent}\ is received, the first message \Emph{OR} the second message above should be displayed randomly. In opposite to the previous exercice, only one message is displayed when an event of type \code{MyEvent}\ is received. For helping you, see the documentation related to the ``event guards''.
	\end{itemize}
\item Create an agent with the type name \code{Agent14}\ that has the following behavior:
	\begin{itemize}
	\item Spawn an agent of type \code{Agent13}.
	\item Every 5 seconds, emit an event of type \code{MyEvent}. The agent of type \code{Agent13} should display \Emph{one} message.
	\end{itemize}
\end{enumerate}


\subsection{Exercice 8 - Behavior}

The goal of this exercice is to introduce the concept of \code{behavior}\ that is supported by a specific keyword in the SARL programming language, and represents a code (or behavior) that is run by an agent. An agent could have multiple behaviors that are run in parallel.

\begin{enumerate}[a)]
\item Create a behavior with the type name \code{Behavior1}\ that has the following behavior:
	\begin{itemize}
	\item Every 1 minute, compute randomly an list of integer values, and save it using the JSON format into the file ``data\_int.json''.
	\end{itemize}
\item Create a behavior with the type name \code{Behavior2}\ that has the following behavior:
	\begin{itemize}
	\item Every 1 minute, compute randomly an list of string values, and save it using the JSON format into the file ``data\_str.json''.
	\end{itemize}
\item Create an agent with the type name \code{Agent15}\ that has the following behavior:
	\begin{itemize}
	\item Create and register a behavior of type \code{Behavior1}\ into the agent.
	\item Create and register a behavior of type \code{Behavior2}\ into the agent.
	\item Run the agent's behaviors.
	\end{itemize}
\end{enumerate}


\end{document}
